% Appendices (after \appendix in main.tex): A = evidence, rules and reference values; B = system, reproducibility, tests.
\newcolumntype{L}{>{\raggedright\arraybackslash}X}

\chapter{Evidence, Rules and Reference Values}
\label{app:evidence}

This appendix lists the evidence, rules and reference values that the
Methodology chapter summarises. Every value is taken from the project file
named in each section.

\section{Intervention Evidence Table}
\label{app:a_evidence}

Every nutritional priority, recipe criterion and scenario claim in HepatoTwin
cites rows of this table (\nolinkurl{data/evidence/intervention_evidence.csv};
Tables~\ref{tab:a_evidence1} and~\ref{tab:a_evidence2}). Null and conflicting
findings are kept so that the twin does not overstate any effect;
observational associations are recorded but never used as effects.

\begin{table}[htbp]
\centering
\footnotesize
\caption{Intervention evidence table, part 1 (weight loss, dietary patterns and exercise).}
\label{tab:a_evidence1}
\begin{tabularx}{\textwidth}{>{\raggedright\arraybackslash}p{2.3cm}LLLp{1.95cm}L}
\toprule
ID & Intervention & Outcome & Design and size & Strength & Source \\
\midrule
VG2015-WL5 & Weight loss $\geq$ 5\% & NASH resolution, NAS & Lifestyle intervention, paired biopsies; 261 of 293 & STRONG & Vilar-Gomez et al. 2015 \\
VG2015-WL7-10 & Weight loss 7--10\% & Liver inflammation (NAS) & Same cohort, 7--10\% subgroup & STRONG & Vilar-Gomez et al. 2015 \\
VG2015-WL10-NASH & Weight loss $\geq$ 10\% & NASH resolution (90\%) & Same cohort, $\geq$ 10\% subgroup & STRONG & Vilar-Gomez et al. 2015 \\
VG2015-WL10-FIB & Weight loss $\geq$ 10\% & Fibrosis regression (45\%) & Same cohort, $\geq$ 10\% subgroup & STRONG & Vilar-Gomez et al. 2015 \\
SAN2021-AST-GGT & Mediterranean diet & AST and GGT (reduced) & Meta-analysis, 10 RCTs, 705 people & MODERATE & Sangouni et al. 2022 \\
SAN2021-ALT & Mediterranean diet & ALT (no significant effect) & Same meta-analysis & WEAK / INCON\-CLUSIVE & Sangouni et al. 2022 \\
STE2026-MEDDIET & Mediterranean diet & ALT, liver stiffness, MRI-PDFF (reduced) & Meta-analysis, 24 RCTs, 4--9 per outcome & MODERATE & Stern et al. 2026 \\
HAI2022-DIET & Dietary interventions & ALT, AST, FLI, steatosis, stiffness & Meta-analysis, 26 articles, 3,037 people & MODERATE & Haigh et al. 2022 \\
HAI2022-CR & Calorie restriction & ALT, steatosis, stiffness & 9 studies within the same review & MODERATE & Haigh et al. 2022 \\
XIO2024-MD-LFD & Mediterranean vs low-fat diet & Liver enzymes and fat (no difference) & Meta-analysis of RCTs & MODERATE & Xiong et al. 2024 \\
STI2023-EX & Exercise training & $\geq$ 30\% fall in MRI liver fat (OR 3.51) & Meta-analysis, 14 studies, 551 people & MODERATE & Stine et al. 2023 \\
STI2023-EXDOSE & Exercise $\geq$ 750 MET-min/week & $\geq$ 30\% fall in MRI liver fat (OR 3.73) & Same meta-analysis & MODERATE & Stine et al. 2023 \\
EASL2024-DIET & Mediterranean-like diet; limit ultra-processed food; avoid sugary drinks & Liver injury & Clinical practice guideline & GUIDELINE & EASL-EASD-EASO 2024 \\
EASL2024-PA & Physical activity $>$ 150 min/week moderate & Steatosis & Clinical practice guideline & GUIDELINE & EASL-EASD-EASO 2024 \\
\bottomrule
\end{tabularx}
\end{table}

\begin{table}[htbp]
\centering
\footnotesize
\caption{Intervention evidence table, part 2 (nutrients, foods and other interventions).}
\label{tab:a_evidence2}
\begin{tabularx}{\textwidth}{>{\raggedright\arraybackslash}p{2.3cm}LLLp{1.95cm}L}
\toprule
ID & Intervention & Outcome & Design and size & Strength & Source \\
\midrule
LUU2018-SAT & Overfeeding 1,000 kcal/day: saturated vs unsaturated fat vs sugars & Liver fat +55\% / +15\% / +33\% & RCT, 38 people, 3 weeks & MODERATE & Luukkonen et al. 2018 \\
CHI2014-HYPER & Fructose adding surplus energy & Liver fat and ALT (raised) & Meta-analysis, 6 hypercaloric trials & MODERATE & Chiu et al. 2014 \\
CHI2014-ISO & Fructose replacing other carbohydrate & Liver fat and ALT (no effect) & Meta-analysis, 7 isocaloric trials & WEAK / INCON\-CLUSIVE & Chiu et al. 2014 \\
LIU2015-COF-FIB & Coffee drinking & Advanced fibrosis (OR 0.73), cirrhosis (OR 0.61) & Observational meta-analysis, 16 studies & OBSERVA\-TIONAL & Liu et al. 2015 \\
EBA2021-COF-FIB & Coffee, highest vs lowest & Significant fibrosis in NAFLD (RR 0.65) & Observational meta-analysis, 5 studies & OBSERVA\-TIONAL & Ebadi et al. 2021 \\
EBA2021-COF-INC & Coffee, highest vs lowest & NAFLD incidence and prevalence (no association) & Observational meta-analysis, 11 studies & OBSERVA\-TIONAL & Ebadi et al. 2021 \\
STE2026-FAST & Fasting & ALT and liver stiffness (reduced) & Meta-analysis, 4--6 RCTs per outcome & MODERATE & Stern et al. 2026 \\
STE2026-LCHF & Low-carbohydrate high-fat or ketogenic diet & ALT (no significant effect) & Meta-analysis, 6 RCTs & WEAK / INCON\-CLUSIVE & Stern et al. 2026 \\
STE2026-OMEGA3 & Omega-3 supplements & ALT (no significant effect) & Meta-analysis, 4 RCTs & WEAK / INCON\-CLUSIVE & Stern et al. 2026 \\
HAL2019-UPF & Ultra-processed vs unprocessed diet & Energy intake (+508 kcal/day), weight & Inpatient crossover RCT, 20 people & MODERATE & Hall et al. 2019 \\
BAW2017-GI & High vs low glycaemic-index diet, 7 days & Liver fat (raised) & Crossover RCT, 8 healthy men & WEAK / INCON\-CLUSIVE & Bawden et al. 2017 \\
\bottomrule
\end{tabularx}
\end{table}

\section{Causal Knowledge Graph}
\label{app:a_graph}

The graph (\nolinkurl{data/evidence/causal_graph.csv}; Table~\ref{tab:a_graph})
links each dietary lever to liver fat. Recipe explanations show the path, for
example ``energy intake $\rightarrow$ body weight $\rightarrow$ liver fat''.
The graph has no cycles and every link cites a source.

\begin{table}[htbp]
\centering
\footnotesize
\caption{Links of the causal knowledge graph.}
\label{tab:a_graph}
\begin{tabularx}{\textwidth}{p{0.8cm}>{\raggedright\arraybackslash}p{4.4cm}>{\raggedright\arraybackslash}p{2.0cm}LL}
\toprule
Link & From $\rightarrow$ to & Relation & Evidence & Note \\
\midrule
CG01 & energy intake $\rightarrow$ body weight & causes (+) & Hall et al. 2011 & About 1 kg per 100 kJ/day \\
CG02 & body weight $\rightarrow$ liver fat & causes (+) & VG2015-WL5, VG2015-WL7-10, HAI2022-CR & $\geq$ 5\% weight loss reduces steatosis \\
CG03 & body weight $\rightarrow$ insulin resistance & causes (+) & Appuhamy et al. 2014 & Weight loss lowers fasting insulin and glucose \\
CG04 & insulin resistance $\rightarrow$ liver fat & causes (+) & Rinella et al. 2023 & Key metabolic driver of MASLD \\
CG05 & saturated fat intake $\rightarrow$ liver fat & causes (+) & LUU2018-SAT, EASL2024-DIET & Raised liver fat more than unsaturated fat or sugars \\
CG06 & added sugar intake $\rightarrow$ liver fat & causes (+) in energy surplus & CHI2014-HYPER, CHI2014-ISO & Only when sugar adds surplus energy \\
CG07 & diet quality $\rightarrow$ liver fat & reduces ($-$) & STE2026-MEDDIET, HAI2022-DIET, EASL2024-DIET & Mediterranean-style diet reduced MRI-PDFF \\
CG08 & physical activity $\rightarrow$ liver fat & reduces ($-$) & STI2023-EX, STI2023-EXDOSE, EASL2024-PA & Works even without large weight loss \\
CG09 & alcohol intake $\rightarrow$ liver fat & causes (+) & Rinella et al. 2023 & Defines MetALD and ALD \\
CG10 & liver fat $\rightarrow$ liver injury & causes (+) & SAN2021-AST-GGT & Enzymes fall as liver fat falls \\
CG11 & liver injury $\rightarrow$ fibrosis & causes (+) & VG2015-WL10-FIB & 10\% weight loss: fibrosis regression in 45\% \\
CG12 & liver fat $\rightarrow$ CAP & is measured by & Wang et al. 2020 & About 5 dB/m per point of MRI-PDFF \\
CG13 & liver injury $\rightarrow$ ALT & is measured by & Koutoukidis et al. 2019 & Marker of injury \\
CG14 & fibrosis $\rightarrow$ liver stiffness & is measured by & EASL-EASD-EASO 2024; Sterling et al. 2006 & Stiffness and FIB-4 estimate fibrosis \\
CG15 & ultra-processed intake $\rightarrow$ energy intake & causes (+) & HAL2019-UPF, EASL2024-DIET & About 500 kcal/day more \\
CG16 & glycaemic load $\rightarrow$ liver fat & causes (+) & BAW2017-GI & Small trial in healthy men \\
\bottomrule
\end{tabularx}
\end{table}

\section{Patient State, Fibrosis Check, Liver Label and Safety Notices}
\label{app:a_state}

The evidence rules fire on yes/no flags derived from the Layer~2
probabilities and routine measurements
(\nolinkurl{src/recipe_scoring/patient_state.py}; Table~\ref{tab:a_flags}).
Flags decide which rules apply; they are not diagnoses.

\begin{table}[htbp]
\centering
\footnotesize
\caption{Patient-state flags and their thresholds.}
\label{tab:a_flags}
\begin{tabularx}{\textwidth}{>{\raggedright\arraybackslash}p{3.4cm}L>{\raggedright\arraybackslash}p{3.6cm}}
\toprule
Flag & Definition & Source \\
\midrule
Steatosis likely & P(CAP $\geq$ 248) $\geq$ 0.432 (F1-selected threshold) & Layer~2 model \\
Moderate--severe steatosis likely & P(CAP $\geq$ 280) $\geq$ 0.365 & Layer~2 model \\
Overweight / obesity & BMI $\geq$ 25 / $\geq$ 30~kg/m$^2$ ($\geq$ 23 / $\geq$ 27.5 for Asian participants) & WHO; WHO Expert Consultation 2004 \\
Central obesity & Waist $>$ 102~cm (men), $>$ 88~cm (women) & NCEP ATP III \\
Elevated glycaemia & HbA1c $\geq$ 5.7\% & ADA prediabetes threshold \\
Dyslipidaemia & Triglycerides $\geq$ 150~mg/dL fasting or $\geq$ 175 non-fasting, or HDL $<$ 40 (men) / $<$ 50 (women)~mg/dL & NCEP ATP III; EAS/EFLM 2016 \\
Recall reliable & Day-1 24-hour recall complete (DR1DRSTZ = 1), a usual day (DR1\_300 = 2) and with plausible energy & NHANES dietary recall \\
\bottomrule
\end{tabularx}
\end{table}

\paragraph{Fibrosis check} (\nolinkurl{src/fibrosis.py}). FIB-4 below 1.3 is low
risk, 1.3--2.67 indeterminate and above 2.67 high (Sterling et al.\ 2006;
EASL-EASD-EASO 2024). No FIB-4 is given under age 18, and a note is added
under age 35 (AASLD 2023). A reliable FibroScan liver stiffness replaces FIB-4
when available: below 8~kPa low, 8 to below 12~kPa indeterminate, 12~kPa or
more high. A high result asks for clinical assessment before any recipes are
shown.

\paragraph{Liver label} (\nolinkurl{src/liver_labels.py}, Rinella et al.\ 2023
nomenclature). Steatosis with at least one cardiometabolic risk factor is
MASLD; with 140--350~g of alcohol a week (women) or 210--420~g (men) it is
MetALD; above that range, ALD; with no risk factor, cryptogenic. Labels support
explanations and are not diagnoses.

\paragraph{Safety notices.} Patients under 18, drinking in the MetALD range or
above, taking glucose-lowering medicine, or pregnant are told to discuss diet
changes with a clinician, and any plan they choose waits for a clinician's
approval. A high fibrosis result also requires approval and withholds
recipes.

\section{Nutritional Priority Rules}
\label{app:a_rules}

Each rule (\nolinkurl{src/recipe_scoring/evidence_rules.py};
Table~\ref{tab:a_rules}) gives a priority with a base level, raised when a
condition holds. When the 24-hour recall is reliable, the level moves up one
step if the patient's own intake misses the same reference value used for
recipes, and down one step if it meets it. No priority carries a calorie or
gram target for the patient: the levels rank which objective matters more, and
every one is marked as needing clinical review.

\begin{table}[htbp]
\centering
\footnotesize
\caption{Nutritional priority rules.}
\label{tab:a_rules}
\begin{tabularx}{\textwidth}{>{\raggedright\arraybackslash}p{2.6cm}>{\raggedright\arraybackslash}p{2.6cm}p{1.2cm}LL>{\raggedright\arraybackslash}p{2.6cm}}
\toprule
Priority & Applies when & Base level & Raised to & Intake check & Evidence \\
\midrule
Evidence-informed weight management & Steatosis likely, with overweight or central obesity & Medium & High with moderate--severe steatosis likely or obesity & None & VG2015-WL5, VG2015-WL7-10, VG2015-WL10-NASH, HAI2022-CR \\
Overall dietary quality & Steatosis likely & Medium & High with moderate--severe steatosis likely & Fibre $\geq$ 14~g per 1,000~kcal & EASL2024-DIET, STE2026-MEDDIET, HAI2022-DIET, SAN2021-AST-GGT, SAN2021-ALT, XIO2024-MD-LFD \\
Recipes lower in saturated fat & Steatosis likely & Medium & High with dyslipidaemia & Saturated fat $\leq$ 10\% of energy & EASL2024-DIET, LUU2018-SAT \\
Recipes lower in added sugars & Steatosis likely & Low & Medium with elevated glycaemia or excess weight & Added sugars $\leq$ 10\% of energy & EASL2024-DIET, CHI2014-HYPER, CHI2014-ISO \\
Less processed ingredients & Steatosis likely & Low & -- & None & EASL2024-DIET, HAL2019-UPF \\
Lower glycaemic load & Steatosis likely and elevated glycaemia & Low & -- & None & BAW2017-GI \\
\bottomrule
\end{tabularx}
\end{table}

\section{Recipe Criteria and Ranking Weights}
\label{app:a_criteria}

Each priority maps to one recipe measure with a reference value from public
guidance (\nolinkurl{src/recipe_scoring/recipe_criteria.py};
Table~\ref{tab:a_criteria}). A criterion scores 100 at or past the reference
and falls linearly to 0 over the tolerance. Gating criteria decide whether a
recipe meets every reference; the two ingredient-derived measures shape the
score but do not gate, because their evidence is weaker.

\begin{table}[htbp]
\centering
\footnotesize
\caption{Recipe criteria, reference values and tolerances.}
\label{tab:a_criteria}
\begin{tabularx}{\textwidth}{>{\raggedright\arraybackslash}p{2.4cm}>{\raggedright\arraybackslash}p{3.0cm}p{1.6cm}p{1.4cm}p{0.8cm}L}
\toprule
Priority & Recipe measure & Reference & Tolerance & Gate & Source \\
\midrule
Weight management & Energy per serving & $\leq$ 600~kcal & 400~kcal & Yes & Public Health England 400-600-600, 2018 \\
Saturated fat & Saturated fat, \% of energy & $\leq$ 10\% & 10 points & Yes & WHO 2023; DGA 2020--2025 \\
Added sugars & Total sugars, \% of energy & $\leq$ 10\% & 15 points & Yes & WHO 2015 (RecipeDB reports total sugars) \\
Dietary quality & Fibre per 1,000~kcal & $\geq$ 14~g & 14~g & Yes & DGA 2020--2025 \\
Less processed & FPro (0--1) & $\leq$ 0.5 & 0.5 & No & Design choice; EASL-EASD-EASO 2024 \\
Glycaemic load & Glycaemic load per serving & $\leq$ 20 & 20 & No & University of Sydney classification \\
\bottomrule
\end{tabularx}
\end{table}

Criterion weights follow the priority level: low 1, medium 2, high 3. The
overall score combines health (weight 0.6), cuisine match (0.2), feasibility
(0.2), acceptability from shrunk RecipeDB ratings (0.1) and personal fit
(0.4); personal fit applies only to recipes that already pass every health
gate (\nolinkurl{src/recipe_scoring/multi_criteria.py}).

\section{Plan Paces}
\label{app:a_paces}

The three paces offered to a patient come from trial prescriptions
(\nolinkurl{src/care_plan.py}, \nolinkurl{data/evidence/simulator_parameters.csv};
Table~\ref{tab:a_paces}). The calorie change and walking together give the
daily energy change; walking is costed at 3.8~MET (2024 Adult Compendium of
Physical Activities).

\begin{table}[htbp]
\centering
\footnotesize
\caption{The three plan paces.}
\label{tab:a_paces}
\begin{tabularx}{\textwidth}{>{\raggedright\arraybackslash}p{2.6cm}L>{\raggedright\arraybackslash}p{2.6cm}L}
\toprule
Pace & Calorie change & Walking & Source \\
\midrule
Gentle & Half the DPP deficit for the weight band & 90 min/week & Derived from DPP \\
Standard (recommended) & DPP deficit, 500--1,000~kcal/day by weight band & 150 min/week & DPP Research Group 2002 \\
Ambitious & Look AHEAD goal: 1,200--1,500~kcal/day ($<$ 114~kg) or 1,500--1,800 ($\geq$ 114~kg) & 175 min/week by month 6 & Wadden et al.\ 2006; Neiberg et al.\ 2012 \\
\bottomrule
\end{tabularx}
\end{table}

% ======================================================================
\chapter{System, Reproducibility and Tests}
\label{app:system}

\section{Architecture}
\label{app:b_architecture}

\begin{enumerate}
    \item \textbf{Layers 2--5} run as Python modules: the liver-state model,
          the evidence engine, recipe ranking and the scenario engine
          (\nolinkurl{src/twin_pipeline.py} runs them together).
    \item \textbf{Backend:} FastAPI with an SQLite database
          (\nolinkurl{src/api/app.py}, \nolinkurl{src/api/db.py}). It stores each
          patient's dated visits, plans and twin runs, so that later visits can
          be compared with what the twin expected.
    \item \textbf{Dashboard:} React and TypeScript, built with Vite
          (\nolinkurl{frontend/}). It calls the backend through \nolinkurl{/api};
          the 3D liver uses three.js.
    \item \textbf{External resources:} RecipeDB for recipes and photo links,
          and optionally an SMTP mail server to email sign-in codes.
\end{enumerate}

\section{API}
\label{app:b_api}

Table~\ref{tab:b_api} lists the backend's endpoints. Patients only ever see
their own record, and never the CAP estimate or grade; clinicians see the
patients of their clinic.

\begin{table}[htbp]
\centering
\footnotesize
\caption{Backend endpoints.}
\label{tab:b_api}
\begin{tabularx}{\textwidth}{>{\raggedright\arraybackslash}p{5.4cm}>{\raggedright\arraybackslash}p{2.4cm}L}
\toprule
Endpoint & Who & What it does \\
\midrule
\texttt{POST /auth/code}, \texttt{POST /auth/session} & Anyone & One-time sign-in code by email, then a session token \\
\texttt{POST /patients} & Patient, clinician & Sign-up: measurements, cuisine, favourite dishes, food rules, clinic code \\
\texttt{GET /patients}, \texttt{GET /patients/\{id\}} & Both & A clinician's patients sorted by alerts; one patient's record \\
\texttt{POST /patients/\{id\}/visits} & Both & A visit or check-in; returns what changed (verdict, chances, tip) \\
\texttt{GET}, \texttt{POST /patients/\{id\}/twin} & Both & The saved twin run for the latest visit, or a new run (patient or clinician view) \\
\texttt{GET .../plan/suggestion}, \texttt{PUT .../plan}, \texttt{POST .../plan/approve} & Both; approval: clinician & Three paces; the chosen plan; clinician approval \\
\texttt{PUT .../plan/style} & Both & Recipe style, step 0 (close to usual food) to 3 (healthiest) \\
\texttt{GET .../today}, \nolinkurl{.../week}, \nolinkurl{.../week/shopping} & Both & Today's next step and meals, the seven-day plan, the shopping list \\
\texttt{GET .../timeline} & Both & Visits, the re-assessment loop and 3-, 6- and 12-month projections \\
\texttt{GET .../clinical} & Clinician & The patient page: CAP estimate and grade, SHAP drivers, fibrosis, priorities, scenarios, explanations \\
\texttt{GET /alerts}, \nolinkurl{/recipes}, \nolinkurl{/dishes}, \nolinkurl{/food/questions} & Mixed & Clinic alerts; recipe details and photos; favourite-dish choices; sign-up food questions \\
\bottomrule
\end{tabularx}
\end{table}

\section{Stored Data}
\label{app:b_data}

\begin{table}[htbp]
\centering
\footnotesize
\caption{Database tables.}
\label{tab:b_data}
\begin{tabularx}{\textwidth}{>{\raggedright\arraybackslash}p{4.4cm}L}
\toprule
Table & Holds \\
\midrule
\texttt{clinics} & Name and joining code \\
\texttt{users}, \texttt{login\_codes}, \texttt{sessions} & Email and role (patient or clinician); hashed one-time codes and session tokens \\
\texttt{patients} & Reference label, cuisine, favourite dishes, food rules \\
\texttt{visits} & Dated NHANES-coded measurements and the check-in answers (meals and walking followed) \\
\texttt{plans} & Start date, pace, calorie change, walking, recipe style, next check-in, status, approver and note \\
\texttt{twin\_runs} & The full twin result for a visit \\
\bottomrule
\end{tabularx}
\end{table}

Table~\ref{tab:b_data} lists the database tables. No names are stored:
patients are known by a label such as P0001 and by their email address.

\section{Dashboard Screens}
\label{app:b_screens}

\begin{table}[htbp]
\centering
\footnotesize
\caption{Dashboard screens.}
\label{tab:b_screens}
\begin{tabularx}{\textwidth}{>{\raggedright\arraybackslash}p{3.8cm}L}
\toprule
Screen & Contents \\
\midrule
Sign-in and sign-up & Emailed code; eight sign-up screens converting answers to NHANES units \\
Today & Next step, body measures, goals, the predict--measure--adjust band, liver summary, today's meals with photos \\
My liver & A 3D liver today beside the liver at 3, 6 or 12 months, coloured by the chance of a lot of fat, then the full results \\
Plan and recipes & Pace choice, waiting for approval, Scenario A or B with dishes and a comparison with usual intake \\
Meals, My plan, Progress, Check-in & The week and shopping list; the plan and recipe-style ladder; weight against the expected range; a four-question check-in \\
Clinician & Patient list by alerts, plan review with a note, and the full patient page \\
\bottomrule
\end{tabularx}
\end{table}

Table~\ref{tab:b_screens} lists the screens; Figures~\ref{fig:m_today},
\ref{fig:m_liver} and~\ref{fig:m_clinician} show three of them.

\section{Reproducing the Results}
\label{app:b_reproduce}

One command rebuilds every result from the raw data:
\texttt{python -m src.rebuild\_all} (\texttt{--list} shows the stages,
\texttt{--from <stage>} resumes). The stages run in this order:

\begin{enumerate}
    \item \textbf{data:} download NHANES 2017--2023 and build the combined
          dataset and the habitual-diet table; collect and profile the
          RecipeDB recipes.
    \item \textbf{model:} train, validate and calibrate the liver-state model;
          conformal intervals; FIB-4 validation.
    \item \textbf{ranking:} evaluate the recipe ranking on 300 held-out
          participants, with and without the 24-hour recall.
    \item \textbf{calibrate, simulate:} calibrate the virtual-patient
          simulator against the trials; run the in-silico evaluation.
    \item \textbf{nhanes3:} download NHANES~III and run the external and
          mortality validation.
    \item \textbf{audits:} the fairness audit and the summary evaluation.
    \item \textbf{figures, tests:} the report figures, then the full test
          suite.
\end{enumerate}

\section{Tests}
\label{app:b_tests}

The Python test suite (\texttt{python -m pytest -q}) covers the model inputs,
evidence rules, ranking, scenarios, re-assessment, safety checks, summaries
and the API, including that patients cannot open the clinician page. Browser
tests in \nolinkurl{frontend/e2e/} drive a headless browser through plan approval,
the patient journey, the 3D liver and the clinician page.

\section{Demo Accounts}
\label{app:b_demo}

The demonstration (\texttt{python -m src.api.seed\_demo}; \texttt{--reset}
restores it) uses a clinic, one clinician and five held-out NHANES
participants who were never used for training: \texttt{dr.rao@example.org}
(clinician); \texttt{demo.patient@example.org} (on diabetes medicine, so his
plan needs approval); and \texttt{patient2}--\texttt{patient5@example.org}
(Mexican cuisine; European vegetarian; American with a 45-minute cooking
limit; Asian pescatarian with little sign of liver fat). Running instructions
are in \nolinkurl{docs/RUN_AND_DEMO.md}.
